%% Relatorio - Reproducao do artigo OptIForest (IJCAI 2023)
%% Escrito sobre o template IEEEtran [conference] exigido pela disciplina.
%% Imagens: subir os 6 arquivos corr_{roc,pr}_{samples,features,anomaly}.png
%% (gerados em results2/analysis/plots/) na raiz do projeto Overleaf.
\documentclass[conference]{IEEEtran}

\usepackage[utf8]{inputenc}
\usepackage[T1]{fontenc}
\usepackage[brazil]{babel}
\usepackage{graphicx}
\usepackage{amsmath,amssymb}
\usepackage{booktabs}
\usepackage{url}
\usepackage{float} % especificador [H]: figura exatamente no ponto do texto
\renewcommand\refname{Refer\^encias}

\usepackage{algorithm}
\usepackage{algpseudocode}
% Traduz os rotulos padrao (em ingles) do algpseudocode.
\renewcommand{\algorithmicwhile}{\textbf{enquanto}}
\renewcommand{\algorithmicdo}{\textbf{fa\c{c}a}}
\renewcommand{\algorithmicend}{\textbf{fim}}
\renewcommand{\algorithmicreturn}{\textbf{retornar}}
\renewcommand{\algorithmicif}{\textbf{se}}
\renewcommand{\algorithmicthen}{\textbf{ent\~ao}}
\renewcommand{\algorithmicelse}{\textbf{sen\~ao}}
\renewcommand{\algorithmicfor}{\textbf{para}}
\renewcommand{\algorithmicrequire}{\textbf{Entrada:}}
\renewcommand{\algorithmicensure}{\textbf{Sa\'ida:}}
\floatname{algorithm}{Algoritmo}

\hyphenation{op-tical net-works semi-conduc-tor}

\begin{document}

\title{Reprodu\c{c}\~ao e Avalia\c{c}\~ao de Generaliza\c{c}\~ao do OptIForest:\\
Detec\c{c}\~ao de Anomalias com Floresta de Isolamento \'Otima}

\author{%
\IEEEauthorblockN{}% mantem o espacamento padrao do bloco de autores do IEEEtran
\centering
\begin{tabular}{c@{\hskip 5em}c}
\shortstack{Pedro Elias\\[2pt]
\textit{Centro de Inform\'atica - UFPE}\\
Recife, Brazil\\
pega@cin.ufpe.br}
&
\shortstack{Tiago Lima\\[2pt]
\textit{Centro de Inform\'atica - UFPE}\\
Recife, Brazil\\
tmsl@cin.ufpe.br}
\\[3ex]
\shortstack{Victor Silva\\[2pt]
\textit{Centro de Inform\'atica - UFPE}\\
Recife, Brazil\\
vsmo@cin.ufpe.br}
&
\shortstack{Vituriano Xisto\\[2pt]
\textit{Centro de Inform\'atica - UFPE}\\
Recife, Brazil\\
vox@cin.ufpe.br}
\end{tabular}
}

\maketitle

\begin{abstract}
A detec\c{c}\~ao n\~ao-supervisionada de anomalias \'e um problema central em
seguran\c{c}a, fraude e diagn\'ostico, no qual as florestas de isolamento
(\emph{Isolation Forest}) se destacam por efici\^encia. Por\'em, a maioria dessas
florestas usa estrutura bin\'aria sem fundamenta\c{c}\~ao te\'orica sobre o fator
de ramifica\c{c}\~ao \'otimo. O artigo de refer\^encia, OptIForest, formaliza
o conceito de efici\^encia de isolamento e mostra que o fator de ramifica\c{c}\~ao
\'otimo \'e $e$, propondo uma floresta que combina \emph{learning-to-hash} por
\emph{clustering}. Neste trabalho, o OptIForest \'e reproduzido sobre os
\emph{benchmarks} do artigo e sua generaliza\c{c}\~ao \'e avaliada em dois conjuntos
de dados externos (\emph{annthyroid} e \emph{fraud}). Reproduz-se a tend\^encia
geral de desempenho do artigo e mostra-se, por an\'alise de correla\c{c}\~ao, que o
AUC-ROC degrada com a dimensionalidade e \'e robusto \`a escala, enquanto o AUC-PR
\'e governado pela raridade da anomalia. Os resultados confirmam que o m\'etodo
generaliza bem em dados separ\'aveis (\emph{fraud}, AUC-ROC $0{,}95$) e exp\~oe seus
limites em anomalias sutis de baixa dimens\~ao (\emph{annthyroid}, AUC-ROC $0{,}57$).
\end{abstract}

\IEEEpeerreviewmaketitle

\section{Introdu\c{c}\~ao}

\textbf{Problema.} Detec\c{c}\~ao de anomalias consiste em identificar inst\^ancias
raras que se desviam do padr\~ao dos dados, tipicamente sem r\'otulos
(n\~ao-supervisionada). \'E interessante por aparecer em dom\'inios cr\'iticos --
fraude financeira, intrus\~ao em redes, diagn\'ostico m\'edico -- nos quais
anomalias s\~ao escassas e caras de rotular.

\textbf{Por que o artigo foi necess\'ario.} A \emph{Isolation Forest} (iForest)
isola anomalias por parti\c{c}\~oes aleat\'orias e \'e eficiente, mas usa \'arvores
bin\'arias por conven\c{c}\~ao. Trabalhos como a LSHiForest mostraram empiricamente
que estruturas multi-ramo podem detectar melhor, por\'em \emph{sem} teoria que
determinasse o fator de ramifica\c{c}\~ao \'otimo. Faltava resposta te\'orica e
pr\'atica para a estrutura ideal de uma floresta de isolamento.

\textbf{Solu\c{c}\~oes existentes e desvantagens.} M\'etodos de isolamento
(iForest, LSHiForest) n\~ao aprendem com os dados; m\'etodos de aprendizado
profundo (REPEN, RDP) aprendem, mas t\^em alto custo computacional; m\'etodos
estat\'isticos (ECOD) s\~ao r\'apidos, por\'em fr\'ageis quando a distribui\c{c}\~ao
foge de suas premissas.

\textbf{Proposta do artigo.} O OptIForest formaliza a \emph{efici\^encia de
isolamento}, prova que o fator de ramifica\c{c}\~ao \'otimo \'e $e \approx 2{,}718$ e
constr\'oi uma floresta que aproxima esse \'otimo combinando \emph{learning-to-hash}
por \emph{clustering}, equilibrando custo e desempenho.

As principais contribui\c{c}\~oes deste trabalho s\~ao:
\begin{itemize}
  \item Reprodu\c{c}\~ao do OptIForest sobre os \emph{benchmarks} do artigo, com
        documenta\c{c}\~ao das diverg\^encias entre os arquivos p\'ublicos e o ap\^endice;
  \item Avalia\c{c}\~ao da generaliza\c{c}\~ao em dois conjuntos externos
        (\emph{annthyroid}, m\'edico; \emph{fraud}, financeiro), com justificativa;
  \item An\'alise de correla\c{c}\~ao entre desempenho e caracter\'isticas dos dados,
        explicando \emph{quando} e \emph{por que} o m\'etodo generaliza.
\end{itemize}

\section{Trabalhos Relacionados}

\textbf{iForest}~\cite{iforest} introduz o isolamento por parti\c{c}\~oes
aleat\'orias; \'e r\'apido, mas n\~ao aprende com os dados e usa \'arvore bin\'aria.
\textbf{LSHiForest}~\cite{lshiforest} generaliza o isolamento via fun\c{c}\~oes de
\emph{hash} sens\'iveis \`a localidade e mostra ganho com estruturas multi-ramo,
sem teoria do fator \'otimo. \textbf{ECOD}~\cite{ecod} estima fun\c{c}\~oes de
distribui\c{c}\~ao emp\'iricas, sendo determin\'istico e r\'apido, mas sens\'ivel \`a
forma da distribui\c{c}\~ao. \textbf{REPEN}~\cite{repen} e \textbf{RDP}~\cite{rdp}
aprendem representa\c{c}\~oes para detec\c{c}\~ao baseada em dist\^ancia, com custo
elevado. O OptIForest se posiciona entre isolamento e aprendizado, herdando a
efici\^encia do primeiro e parte da capacidade do segundo. A
Tabela~\ref{tab:related} resume os m\'etodos.

\begin{table}[h]
\centering
\caption{S\'intese dos trabalhos relacionados.}
\label{tab:related}
\footnotesize
\begin{tabular}{llp{2.4cm}}
\toprule
\textbf{Trabalho} & \textbf{Tipo} & \textbf{Limita\c{c}\~ao} \\
\midrule
iForest~\cite{iforest}      & Isolamento aleat\'orio & N\~ao aprende; \'arvore bin\'aria \\
LSHiForest~\cite{lshiforest}& Isolamento via LSH     & Sem teoria do fator \'otimo \\
ECOD~\cite{ecod}            & Estat\'istico          & Sens\'ivel \`a distribui\c{c}\~ao \\
REPEN~\cite{repen}          & Aprendizado profundo   & Custo computacional alto \\
RDP~\cite{rdp}              & Aprendizado profundo   & Custo computacional alto \\
OptIForest~\cite{optiforest}& Isolamento + \emph{learning} & Sens\'ivel ao limiar $\epsilon$ \\
\bottomrule
\end{tabular}
\end{table}

\section{Modelo de Ameaça}
Embora o OptIForest seja um detector de anomalias de prop\'osito geral, no contexto
desta disciplina ele atua como n\'ucleo de um \emph{sistema de detec\c{c}\~ao de
intrus\~ao} (IDS) baseado em anomalia: parte dos \emph{benchmarks} reproduzidos s\~ao,
de fato, tra\c{c}os de seguran\c{c}a -- \emph{http} e \emph{kddcup99} (conex\~oes de
rede do KDDCup'99), \emph{smtp} (tr\'afego de \emph{e-mail}) e \emph{fraud} (transa\c{c}\~oes
de cart\~ao). Nesses casos, a \textbf{anomalia} a ser ranqueada \'e justamente a
tentativa de intrus\~ao ou fraude.

\textbf{Ataques considerados.} Intrus\~oes de rede (varreduras, \emph{DoS}, acessos
n\~ao-autorizados) e fraude financeira, tratados como inst\^ancias raras cujo vetor de
\emph{features} se desvia estatisticamente do tr\'afego/transa\c{c}\~oes leg\'itimas.
\textbf{Premissas} para o detector funcionar: (i)~ataques s\~ao \emph{raros} frente ao
volume normal (verdadeiro nos dados: $0{,}03\%$ a $\sim 39\%$ de anomalias);
(ii)~ataques s\~ao \emph{separ\'aveis} no espa\c{c}o de \emph{features}, isto \'e, exigem
menos parti\c{c}\~oes para serem isolados; (iii)~o detector \'e treinado em modo
n\~ao-supervisionado sobre dados majoritariamente normais, sem r\'otulos de ataque.
\textbf{Fora do escopo:} n\~ao se modela um \emph{advers\'ario adaptativo} (ex.\
ataques de evas\~ao/envenenamento que manipulem deliberadamente as \emph{features}
para parecer normais); o modelo assume que o atacante n\~ao tem conhecimento do
detector. Esse limite \'e retomado em Trabalhos Futuros.

\section{Sistema Proposto pelo Artigo de Refer\^encia}

\textbf{Vis\~ao geral.} O OptIForest \'e um \emph{ensemble} de \'arvores de
isolamento \'otimas. \textbf{Entrada:} um conjunto de dados $D$ sem r\'otulos.
\textbf{Sa\'ida:} um \emph{score} de anomalia por inst\^ancia (quanto menor o
caminho de isolamento, mais an\^omala). Cada \'arvore tem duas partes: camadas
inferiores produzidas por uma floresta LSH (\emph{LSHiForest}) e camadas superiores
refinadas por \emph{clustering} hier\'arquico aglomerativo -- a parte que ``aprende''
com os dados.

\textbf{Componentes (entrada $\rightarrow$ fun\c{c}\~ao $\rightarrow$ sa\'ida).}
(1)~\emph{Amostragem}: de $D$ extrai sub-amostras independentes (uma por \'arvore),
gerando diversidade no \emph{ensemble}. (2)~\emph{LSHTree}: sobre cada sub-amostra,
particiona recursivamente os dados com fun\c{c}\~oes de \emph{hash} sens\'iveis \`a
localidade; produz uma \'arvore de isolamento em que cada n\'o guarda centro,
tamanho e a fun\c{c}\~ao \emph{hash} usada. Para dist\^ancia euclidiana usa-se a fam\'ilia
E2LSH,
\begin{equation}
  h(x) = \left\lfloor \frac{a^{\top}x + b}{w} \right\rfloor,
  \label{eq:e2lsh}
\end{equation}
com vetor aleat\'orio $a$, deslocamento $b$ e largura de \emph{bin} $w$, de modo que
pontos pr\'oximos caiam no mesmo compartimento com alta probabilidade.
(3)~\emph{Corte por $\epsilon$ + clustering hier\'arquico}: refina as camadas
superiores da \'arvore (Alg.~\ref{alg:opt}). (4)~\emph{Floresta}: agrega as
\'arvores \'otimas. (5)~\emph{Score}: dada uma inst\^ancia, devolve seu grau de
anomalia (Eq.~\ref{eq:score}).

\textbf{Efici\^encia de isolamento e fator $e$.} O artigo define a efici\^encia de
isolamento e formula um problema de otimiza\c{c}\~ao cuja solu\c{c}\~ao indica fator de
ramifica\c{c}\~ao \'otimo $v \approx e$. Na pr\'atica, usam-se ramos bin\'arios e
tern\'arios com distribui\c{c}\~ao exponencial para aproximar $e$.

\textbf{Aprendizado por clustering.} A \'arvore LSH \'e cortada por um limiar
$\epsilon$ (\emph{cut threshold}): n\'os com tamanho $\le \epsilon$ tornam-se
\emph{clusters} iniciais, que s\~ao fundidos minimizando a \emph{distor\c{c}\~ao}
\begin{equation}
  \mathrm{dist}(\mathbb{C}) = \sum_{i=1}^{v} \lVert \mu_{\mathcal{C}_i} - \mu_{\mathbb{C}} \rVert \cdot n_i,
  \label{eq:distortion}
\end{equation}
onde $\mu$ s\~ao centros ponderados pelo tamanho $n_i$ dos \emph{clusters}. O
par\^ametro $\epsilon$ controla o \emph{grau de aprendizado}: $\epsilon$ pequeno
aprende mais (mais \emph{clusters}); $\epsilon$ grande degenera na LSHiForest.

\begin{algorithm}[h]
\caption{Constru\c{c}\~ao da \'arvore de isolamento \'otima (resumo).}
\label{alg:opt}
\begin{algorithmic}[1]
\Require conjunto $D$, limiar $\epsilon$
\State Treinar LSHiForest para obter a \'arvore $T_{LSH}$
\State Cortar $T_{LSH}$ pelo limiar $\epsilon$ $\rightarrow$ \emph{clusters} iniciais
\While{n\'umero de \emph{clusters} $> 1$}
  \State Sortear fator de ramifica\c{c}\~ao $v \in \{2,3\}$ (aprox.\ $e$)
  \State Fundir os $v$ \emph{clusters} de menor distor\c{c}\~ao (Eq.~\ref{eq:distortion})
\EndWhile
\State \Return \'arvore \'otima $T_{Opt}$
\end{algorithmic}
\end{algorithm}

\textbf{Score de anomalia.} Como na iForest, o \emph{score} decai com a profundidade
m\'edia de isolamento $h(x)$ na floresta, normalizada pelo comprimento esperado de
caminho $c(n)$ de uma \'arvore com $n$ amostras:
\begin{equation}
  s(x) = 2^{-\,h(x)/c(n)}.
  \label{eq:score}
\end{equation}
Inst\^ancias isoladas com poucas parti\c{c}\~oes (caminho curto) recebem $s(x)$
pr\'oximo de $1$ e s\~ao ranqueadas como mais an\^omalas.

\section{Metodologia}

\textbf{Protocolo.} O OptIForest \'e \textbf{n\~ao-supervisionado}: n\~ao h\'a divis\~ao
treino/valida\c{c}\~ao/teste. O modelo \'e ajustado e avaliado sobre \emph{todo} o
conjunto, e as anomalias servem apenas como gabarito para o c\'alculo das m\'etricas.
Por isso reporta-se a contagem por classe (normal vs.\ anomalia) do conjunto completo,
n\~ao por parti\c{c}\~ao. Cada configura\c{c}\~ao roda $15$ vezes (sementes $42$--$56$),
com $100$ \'arvores e o limiar $\epsilon$ pela regra do artigo ($512$ para
\emph{ad}/\emph{vowel}; $403$ para $>10^4$ amostras ou $>10^3$ \emph{features};
$55$ caso contr\'ario).

\subsection{Dados usados pelo artigo de refer\^encia}
O artigo avalia $20$ \emph{benchmarks} tabulares. Reproduzem-se todos os
dispon\'iveis a partir de fontes p\'ublicas (ADBench~\cite{adbench} e OpenML). Tr\^es
exemplos representativos por diversidade de caracter\'isticas:
\begin{itemize}
  \item \emph{arrhythmia} (m\'edico, alta dimens\~ao): $452$ amostras, $274$
        \emph{features}; $66$ anomalias / $386$ normais.
  \item \emph{mnist} (imagem, dimens\~ao m\'edia): $7\,603$ amostras, $100$
        \emph{features}; $700$ anomalias / $6\,903$ normais.
  \item \emph{cover} (grande, anomalia rara): $286\,048$ amostras, $10$
        \emph{features}; $2\,746$ anomalias / $283\,302$ normais.
\end{itemize}
A contagem por classe dos $20$ conjuntos deriva diretamente das colunas
\#Amostras e Anom.\% da Tabela~\ref{tab:repro} (\emph{normais} $=$
$\#\text{Amostras}\times(1-\text{Anom.\%})$).

\subsection{Novo conjunto de dados escolhido}
Para avaliar generaliza\c{c}\~ao, escolheram-se dois conjuntos \textbf{fora} dos $20$ do
artigo, ambos do ADBench (mesmo formato, sem c\'odigo adicional de I/O):
\begin{itemize}
  \item \emph{annthyroid} (m\'edico, tireoide): $7\,200$ amostras, $6$
        \emph{features}; $534$ anomalias / $6\,666$ normais. \emph{Por qu\^e:}
        regime de baix\'issima dimens\~ao, contrastando com \emph{arrhythmia}
        (m\'edico, mas $274$ \emph{features}), isolando o efeito da dimensionalidade.
  \item \emph{fraud} (financeiro, cart\~ao de cr\'edito): $284\,807$ amostras, $29$
        \emph{features}; $484$ anomalias / $284\,323$ normais. \emph{Por qu\^e:}
        dom\'inio novo e anomalia ultra-rara ($0{,}17\%$), abaixo do menor caso do
        artigo (\emph{cover}, $0{,}96\%$), testando robustez em desbalanceamento
        extremo.
\end{itemize}

\subsection{M\'etricas de avalia\c{c}\~ao}
Seguindo o artigo, utilizam-se duas m\'etricas baseadas em ranqueamento, ambas
insens\'iveis \`a escolha de um limiar de decis\~ao. A \textbf{AUC-ROC} \'e a \'area
sob a curva que relaciona a taxa de verdadeiros positivos
($\mathrm{TPR}=\frac{TP}{TP+FN}$) \`a de falsos positivos
($\mathrm{FPR}=\frac{FP}{FP+TN}$):
\begin{equation}
  \mathrm{AUC\text{-}ROC} = \int_{0}^{1} \mathrm{TPR}(\mathrm{FPR})\, d(\mathrm{FPR}).
  \label{eq:auroc}
\end{equation}
Mede a capacidade de ranquear anomalias acima de normais e \'e \emph{insens\'ivel \`a
preval\^encia}. A \textbf{AUC-PR} \'e a \'area sob a curva precis\~ao--revoca\c{c}\~ao,
com $\mathrm{Prec}=\frac{TP}{TP+FP}$ e $\mathrm{Rec}=\frac{TP}{TP+FN}$:
\begin{equation}
  \mathrm{AUC\text{-}PR} = \int_{0}^{1} \mathrm{Prec}(\mathrm{Rec})\, d(\mathrm{Rec}).
  \label{eq:aupr}
\end{equation}
\'E \emph{sens\'ivel ao desbalanceamento}: como sua linha de base \'e a pr\'opria
taxa de anomalia, expor o desempenho em cen\'arios de anomalia ultra-rara
(\emph{fraud}, \emph{smtp}) \'e justamente o que se deseja observar.

\subsection{Reprodu\c{c}\~ao, ambiente e otimiza\c{c}\~oes de execu\c{c}\~ao}
Este trabalho parte do c\'odigo p\'ublico do OptIForest~\cite{optiforest}
(\texttt{xiagll/optiforest}), reorganizado em um \emph{runner} que baixa os dados
(ADBench~\cite{adbench} e OpenML), executa as $15$ \emph{runs} por configura\c{c}\~ao e
consolida as m\'etricas; todo o c\'odigo comentado e o hist\'orico de modifica\c{c}\~oes
est\~ao no reposit\'orio GitHub da equipe. Conjuntos com diverg\^encia em rela\c{c}\~ao ao
ap\^endice (\emph{wine}, \emph{waveform}; \emph{kddcup99} sem arquivo \'unico, suprido
pelo \emph{http}) s\~ao sinalizados via uma \emph{flag} \texttt{paper\_shape\_match} na
sa\'ida.

\textbf{Ambiente.} Os conjuntos grandes (\emph{cover}, \emph{donors}, \emph{fraud},
\emph{census}) foram executados no \textbf{cluster Apuana} do CIn-UFPE, via SLURM.
Para evitar \emph{oversubscription} entre processos e \emph{threads} de \'algebra
linear, fixou-se \texttt{OMP\_NUM\_THREADS=1} (e BLAS/MKL) e paralelizou-se no n\'ivel
das tarefas: a \emph{flag} \texttt{--flat-workers} forma um \'unico \emph{pool} sobre
todos os pares (\emph{dataset}, \emph{run}), mantendo os n\'ucleos ocupados
independentemente das fronteiras de \emph{dataset}.

\textbf{Otimiza\c{c}\~ao de implementa\c{c}\~ao.} O custo dominante da fase de
\emph{clustering} estava no m\'etodo \texttt{HierTree.index\_merged\_node}, que
escolhe os $v\in\{2,3\}$ \emph{clusters} de menor distor\c{c}\~ao
(Eq.~\ref{eq:distortion}) com \emph{loops} Python $O(n^2)$/$O(n^3)$. A busca foi
\textbf{vetorizada}: derivando $\lVert \mu_{\text{nc}} - C_k\rVert$ em
fun\c{c}\~ao apenas das dist\^ancias par-a-par, calcula-se a matriz de dist\^ancias
ponderadas com uma \'unica chamada \texttt{pdist}/\texttt{squareform} do SciPy e
toma-se o \texttt{argmin}. A reescrita preserva o mesmo par/trio escolhido (a menos de
arredondamento de ponto flutuante) e, portanto, \textbf{n\~ao altera os resultados}
(mesmo AUC), apenas o tempo de execu\c{c}\~ao -- por isso \'e tratada como detalhe de
implementa\c{c}\~ao, e n\~ao como modifica\c{c}\~ao do m\'etodo. O caso tern\'ario
($v=3$) era cr\'itico: na configura\c{c}\~ao do artigo (\emph{branch}~$=0$, que aproxima
$e$ misturando fus\~oes bin\'arias e tern\'arias), o \emph{loop} $O(n^3)$ original
\textbf{n\~ao terminava} em tempo h\'abil nos conjuntos grandes (\emph{mnist},
\emph{cover}). Vetorizar tamb\'em o $v=3$ -- por uma identidade an\'aloga sobre as
dist\^ancias par-a-par -- foi o que tornou a reprodu\c{c}\~ao desses conjuntos
fact\'ivel dentro do prazo.

\section{Resultados e Discuss\~oes}

\subsection{Resultados da reprodu\c{c}\~ao do artigo de refer\^encia}
A Tabela~\ref{tab:repro} resume os resultados reproduzidos, que seguem a
tend\^encia geral do artigo: desempenho alto em \emph{http}, \emph{shuttle},
\emph{wine} e \emph{ionosphere}, e mais baixo em \emph{census} e \emph{campaign}.
A Tabela~\ref{tab:comp} compara o AUC-ROC reproduzido com o do artigo (valores
conferidos na Tabela~1 do PDF oficial). \textbf{A maioria dos conjuntos reproduz com
alta fidelidade:} doze dos vinte ficam dentro de $\pm 2$ pontos do artigo
(\emph{ad} $+0{,}9$, \emph{donors} $+0{,}1$, \emph{http} $-0{,}0$, \emph{spam}
$-0{,}0$, \emph{census} $-0{,}3$, \emph{ionosphere} $+0{,}6$, \emph{waveform}
$+0{,}7$, \emph{shuttle} $+0{,}9$, \emph{vowel} $+1{,}1$, \emph{celeba} $-1{,}3$,
\emph{backdoor} $-1{,}8$, \emph{kddcup99} $+2{,}0$), o que valida a corre\c{c}\~ao da
implementa\c{c}\~ao -- incluindo a otimiza\c{c}\~ao vetorizada do \emph{clustering}.
As maiores diverg\^encias concentram-se onde o artefato p\'ublico difere do usado no
artigo: \emph{wine} ($+33{,}3$, pois o ADBench usa o conjunto ODDS~\cite{odds}
$129\times13$, n\~ao o $5\,318\times11$ do ap\^endice) e os conjuntos de maior
diverg\^encia (\emph{mnist} $-25{,}3$, \emph{cover} $-22{,}5$, \emph{satellite}
$-12{,}6$), investigados a seguir.

\textbf{Diagn\'ostico das maiores diverg\^encias.} Para \emph{mnist} e \emph{cover},
descartou-se a hip\'otese de dado corrompido: uma \emph{Isolation Forest} cl\'assica
(\emph{scikit-learn}) aplicada ao \emph{mesmo} arquivo \texttt{24\_mnist.npz} recupera
AUC-ROC $\approx 0{,}80$, coerente com a literatura -- ou seja, o conjunto p\'ublico
est\'a correto. Dois ind\'icios apontam, ent\~ao, para uma diferen\c{c}a entre o
c\'odigo p\'ublico e o efetivamente usado no artigo: (i)~a queda persiste apenas na
implementa\c{c}\~ao do OptIForest, n\~ao no isolamento puro sobre o mesmo dado; e
(ii)~os tempos reportados no artigo (\emph{mnist} $\approx 132$\,s) s\~ao cerca de
$20\times$ menores que os obtidos aqui mesmo \emph{ap\'os} a otimiza\c{c}\~ao da fase de
\emph{clustering} ($\approx 2\,400$\,s). A discrep\^ancia de tempo, em particular, \'e
incompat\'ivel com a mesma rotina de \emph{clustering} hier\'arquico e sugere que o
artefato p\'ublico de \texttt{xiagll/optiforest} n\~ao corresponde \`a vers\~ao
experimental que gerou a Tabela~1 -- e n\~ao um erro de reprodu\c{c}\~ao desta equipe.

\textbf{Compara\c{c}\~ao com outros m\'etodos.} No artigo, o OptIForest lidera o
AUC-ROC m\'edio ($83{,}9\%$), \`a frente da LSHiForest ($81{,}9\%$), da ECOD
($78{,}8\%$) e da iForest cl\'assica ($74{,}3\%$)~\cite{optiforest}. A reprodu\c{c}\~ao
deste trabalho recupera a mesma ordena\c{c}\~ao qualitativa entre datasets, confirmando
que o ganho do \emph{learning-to-hash} sobre o isolamento puramente aleat\'orio se
mant\'em fora do ambiente original dos autores.

\begin{table}[t]
\centering
\caption{Resultados reproduzidos do OptIForest (AUC-ROC e AUC-PR em \%,
m\'edia $\pm$ desvio sobre 15 execu\c{c}\~oes). Datasets novos em \textbf{negrito}.}
\label{tab:repro}
\footnotesize
\setlength{\tabcolsep}{3.5pt}
\begin{tabular}{lrrrr}
\toprule
Dataset & \#Amostras & Anom.\% & AUC-ROC & AUC-PR \\
\midrule
ad & 3279 & 14.00 & 78.3 $\pm$ 0.4 & 45.5 $\pm$ 1.6 \\
\textbf{annthyroid} & 7200 & 7.42 & 56.8 $\pm$ 2.0 & 11.8 $\pm$ 1.3 \\
arrhythmia & 452 & 14.60 & 77.4 $\pm$ 1.3 & 53.8 $\pm$ 2.9 \\
backdoor & 95329 & 2.44 & 90.9 $\pm$ 1.1 & 29.4 $\pm$ 6.3 \\
campaign & 41188 & 11.27 & 68.9 $\pm$ 1.3 & 25.6 $\pm$ 1.8 \\
cardio & 1831 & 9.61 & 85.9 $\pm$ 1.4 & 48.4 $\pm$ 2.2 \\
celeba & 202599 & 2.24 & 77.9 $\pm$ 1.7 & 7.7 $\pm$ 1.0 \\
census & 299285 & 6.20 & 67.5 $\pm$ 1.1 & 8.7 $\pm$ 0.3 \\
cover & 286048 & 0.96 & 68.4 $\pm$ 4.0 & 1.5 $\pm$ 0.3 \\
donors & 619326 & 5.93 & 77.2 $\pm$ 4.6 & 12.0 $\pm$ 2.0 \\
\textbf{fraud} & 284807 & 0.17 & 94.8 $\pm$ 0.4 & 41.7 $\pm$ 5.7 \\
http & 567498 & 0.39 & 99.4 $\pm$ 0.0 & 35.8 $\pm$ 0.9 \\
ionosphere & 351 & 35.90 & 94.0 $\pm$ 0.4 & 92.6 $\pm$ 0.4 \\
kddcup99 & 567498 & 0.39 & 99.4 $\pm$ 0.0 & 35.8 $\pm$ 0.9 \\
mnist & 7603 & 9.21 & 60.2 $\pm$ 2.3 & 15.1 $\pm$ 2.8 \\
satellite & 6435 & 31.64 & 66.0 $\pm$ 1.6 & 64.7 $\pm$ 1.7 \\
shuttle & 49097 & 7.15 & 98.9 $\pm$ 0.3 & 82.5 $\pm$ 3.0 \\
smtp & 95156 & 0.03 & 86.2 $\pm$ 1.4 & 60.6 $\pm$ 3.6 \\
spam & 4207 & 39.91 & 71.1 $\pm$ 0.4 & 58.5 $\pm$ 0.7 \\
vowel & 1456 & 3.43 & 91.1 $\pm$ 0.9 & 31.8 $\pm$ 3.2 \\
waveform & 3443 & 2.90 & 75.1 $\pm$ 1.2 & 6.0 $\pm$ 0.5 \\
wine & 129 & 7.75 & 99.7 $\pm$ 0.2 & 96.8 $\pm$ 1.9 \\
\bottomrule
\end{tabular}
\end{table}

\begin{table}[t]
\centering
\caption{Reprodu\c{c}\~ao vs.\ paper (AUC-ROC \%). Apenas datasets do paper. Valores
do paper conferidos na Tabela~1 do artigo oficial~\cite{optiforest}.}
\label{tab:comp}
\footnotesize
\begin{tabular}{lrrr}
\toprule
Dataset & Paper & Reproduzido & $\Delta$ \\
\midrule
ad & 77.4 & 78.3 & +0.9 \\
arrhythmia & 79.6 & 77.4 & -2.2 \\
backdoor & 92.7 & 90.9 & -1.8 \\
campaign & 74.7 & 68.9 & -5.8 \\
cardio & 92.8 & 85.9 & -6.9 \\
celeba & 79.2 & 77.9 & -1.3 \\
census & 67.8 & 67.5 & -0.3 \\
cover & 90.9 & 68.4 & -22.5 \\
donors & 77.1 & 77.2 & +0.1 \\
http & 99.4 & 99.4 & -0.0 \\
ionosphere & 93.4 & 94.0 & +0.6 \\
kddcup99 & 97.4 & 99.4 & +2.0 \\
mnist & 85.5 & 60.2 & -25.3 \\
satellite & 78.6 & 66.0 & -12.6 \\
shuttle & 98.0 & 98.9 & +0.9 \\
smtp & 92.4 & 86.2 & -6.2 \\
spam & 71.1 & 71.1 & -0.0 \\
vowel & 90.0 & 91.1 & +1.1 \\
waveform & 74.4 & 75.1 & +0.7 \\
wine & 66.4 & 99.7 & +33.3 \\
\bottomrule
\end{tabular}
\end{table}

\subsection{Resultados no novo conjunto de dados}
\emph{fraud} obteve AUC-ROC $0{,}948 \pm 0{,}004$ e AUC-PR $0{,}417 \pm 0{,}057$:
mesmo com $0{,}17\%$ de anomalias, as fraudes s\~ao bem separ\'aveis e o m\'etodo
ranqueia com excel\^encia. \emph{annthyroid} obteve AUC-ROC $0{,}568 \pm 0{,}020$ e
AUC-PR $0{,}118 \pm 0{,}013$: abaixo da faixa t\'ipica de iForest, pois a anomalia \'e
sutil e o limiar $\epsilon = 55$ (regra do artigo para $\le 10^4$ amostras) limita o
aprendizado em baixa dimens\~ao. O contraste entre os dois casos delimita
\emph{quando} o m\'etodo generaliza.

\subsection{Correla\c{c}\~ao entre desempenho e caracter\'isticas}
A Tabela~\ref{tab:corr} resume a correla\c{c}\~ao (Pearson e Spearman, que concordam em
sinal) entre desempenho e tr\^es caracter\'isticas dos $22$ conjuntos -- n\'umero de
amostras, n\'umero de \emph{features} e taxa de anomalia. As
Figuras~\ref{fig:corr-roc-samples}--\ref{fig:corr-pr-anomaly} detalham cada rela\c{c}\~ao;
a seguir discute-se uma caracter\'istica por vez.

\begin{table}[t]
\centering
\caption{Correla\c{c}\~ao (Pearson~/~Spearman) entre desempenho e caracter\'isticas.}
\label{tab:corr}
\footnotesize
\setlength{\tabcolsep}{4pt}
\begin{tabular}{llrr}
\toprule
Caracter\'istica & M\'etrica & Pearson & Spearman \\
\midrule
N\textsuperscript{o} amostras   & AUC-ROC & $0{,}03$  & $-0{,}03$ \\
N\textsuperscript{o} amostras   & AUC-PR  & $-0{,}54$ & $-0{,}50$ \\
N\textsuperscript{o} features   & AUC-ROC & $-0{,}35$ & $-0{,}38$ \\
N\textsuperscript{o} features   & AUC-PR  & $-0{,}07$ & $-0{,}05$ \\
Taxa de anomalia                & AUC-ROC & $-0{,}24$ & $-0{,}35$ \\
Taxa de anomalia                & AUC-PR  & $0{,}50$  & $0{,}40$  \\
\bottomrule
\end{tabular}
\end{table}

\textbf{N\'umero de amostras.} O AUC-ROC \textbf{independe da escala}
($r \approx 0$, Fig.~\ref{fig:corr-roc-samples}): o m\'etodo ranqueia igualmente bem de
$129$ (\emph{wine}) a $619$ mil amostras (\emph{donors}), evid\^encia direta de boa
generaliza\c{c}\~ao em tamanho. J\'a o AUC-PR \textbf{cai} com o n\'umero de amostras
($r \approx -0{,}54$, Fig.~\ref{fig:corr-pr-samples}), mas esse efeito \'e
\emph{mec\^anico}: conjuntos maiores tendem a ter menor taxa de anomalia, e a linha de
base do AUC-PR \'e a pr\'opria preval\^encia -- n\~ao se trata de piora real do
ranqueamento.

\begin{figure}[H]
\centering
\includegraphics[width=\linewidth]{corr_roc_samples.png}
\caption{AUC-ROC vs.\ n\'umero de amostras (vermelho = conjuntos novos).}
\label{fig:corr-roc-samples}
\end{figure}

\begin{figure}[H]
\centering
\includegraphics[width=\linewidth]{corr_pr_samples.png}
\caption{AUC-PR vs.\ n\'umero de amostras (vermelho = conjuntos novos).}
\label{fig:corr-pr-samples}
\end{figure}

\textbf{N\'umero de \emph{features}.} O AUC-ROC \textbf{degrada com a dimensionalidade}
($r \approx -0{,}35$, Fig.~\ref{fig:corr-roc-features}), coerente com a maldi\c{c}\~ao da
dimens\~ao sobre as fun\c{c}\~oes de \emph{hash}: em espa\c{c}os de alta dimens\~ao as
dist\^ancias se concentram e o isolamento perde poder discriminativo (\emph{mnist},
\emph{arrhythmia}). O AUC-PR \'e praticamente \textbf{indiferente} \`a dimens\~ao
($r \approx -0{,}07$, Fig.~\ref{fig:corr-pr-features}), pois \'e dominado pela raridade
da anomalia, n\~ao pela geometria do espa\c{c}o.

\begin{figure}[H]
\centering
\includegraphics[width=\linewidth]{corr_roc_features.png}
\caption{AUC-ROC vs.\ n\'umero de \emph{features} (vermelho = conjuntos novos).}
\label{fig:corr-roc-features}
\end{figure}

\begin{figure}[H]
\centering
\includegraphics[width=\linewidth]{corr_pr_features.png}
\caption{AUC-PR vs.\ n\'umero de \emph{features} (vermelho = conjuntos novos).}
\label{fig:corr-pr-features}
\end{figure}

\textbf{Taxa de anomalia.} O AUC-PR \textbf{cresce} com a taxa de anomalia
($r \approx +0{,}50$, Fig.~\ref{fig:corr-pr-anomaly}) -- novamente em parte
\emph{mec\^anico}, j\'a que a preval\^encia \'e a sua linha de base. O AUC-ROC tem
correla\c{c}\~ao levemente \textbf{negativa} ($r \approx -0{,}24$,
Fig.~\ref{fig:corr-roc-anomaly}). Por ser insens\'ivel \`a preval\^encia, o AUC-ROC \'e
o crit\'erio mais justo para comparar conjuntos. Esses padr\~oes explicam diretamente os
dois casos novos: \emph{fraud} (baixa dimens\~ao, separ\'avel) segue o padr\~ao e
generaliza, enquanto \emph{annthyroid} (anomalia sutil) \'e o contraexemplo que revela
os limites do m\'etodo.

\begin{figure}[H]
\centering
\includegraphics[width=\linewidth]{corr_roc_anomaly.png}
\caption{AUC-ROC vs.\ taxa de anomalia (vermelho = conjuntos novos).}
\label{fig:corr-roc-anomaly}
\end{figure}

\begin{figure}[H]
\centering
\includegraphics[width=\linewidth]{corr_pr_anomaly.png}
\caption{AUC-PR vs.\ taxa de anomalia (vermelho = conjuntos novos).}
\label{fig:corr-pr-anomaly}
\end{figure}

\section{Conclus\~oes e Trabalhos Futuros}
Este trabalho reproduziu o OptIForest e confirmou sua tend\^encia de desempenho,
documentando as diverg\^encias decorrentes de arquivos p\'ublicos distintos do
ap\^endice do artigo. A avalia\c{c}\~ao em dois conjuntos externos mostrou que o m\'etodo
\textbf{generaliza bem em escala} e em dados separ\'aveis (\emph{fraud}), mas
\textbf{sofre com alta dimensionalidade} e com anomalias sutis em baixa dimens\~ao
(\emph{annthyroid}). \textbf{Limita\c{c}\~oes:} aus\^encia de divis\~ao
treino/teste (intr\'inseca ao cen\'ario n\~ao-supervisionado), sensibilidade ao
limiar $\epsilon$ fixado por tamanho, e custo de predi\c{c}\~ao elevado em conjuntos
grandes (\emph{cover}). \textbf{Trabalhos futuros:} tornar $\epsilon$ adaptativo \`a
densidade local; investigar fun\c{c}\~oes de \emph{hash} angulares em alta dimens\~ao;
avaliar a \textbf{robustez a advers\'arios adaptativos} (ataques de evas\~ao/
envenenamento), fora do modelo de amea\c{c}a atual; e a vers\~ao federada citada pelo
pr\'oprio artigo.

\begin{thebibliography}{9}
\bibitem{optiforest} H.~Xiang, X.~Zhang, H.~Hu, L.~Qi, W.~Dou, M.~Dras, A.~Beheshti,
  and X.~Xu, ``OptIForest: Optimal Isolation Forest for Anomaly Detection,'' in
  \emph{Proc. 32nd Int. Joint Conf. Artif. Intell. (IJCAI)}, 2023, pp.~4486--4494.
\bibitem{iforest} F.~T.~Liu, K.~M.~Ting, and Z.-H.~Zhou, ``Isolation Forest,'' in
  \emph{Proc. IEEE ICDM}, 2008, pp.~413--422.
\bibitem{lshiforest} X.~Zhang, W.~Dou, Q.~He, R.~Zhou, C.~Leckie, R.~Kotagiri, and
  Z.~Salcic, ``LSHiForest: A Generic Framework for Fast Tree Isolation Based
  Ensemble Anomaly Analysis,'' in \emph{Proc. IEEE ICDE}, 2017, pp.~983--994.
\bibitem{ecod} Z.~Li \emph{et al.}, ``ECOD: Unsupervised Outlier Detection Using
  Empirical Cumulative Distribution Functions,'' \emph{IEEE TKDE}, 2022.
\bibitem{repen} G.~Pang, L.~Cao, L.~Chen, and H.~Liu, ``Learning Representations of
  Ultrahigh-dimensional Data for Random Distance-based Outlier Detection,'' in
  \emph{Proc. ACM KDD}, 2018.
\bibitem{rdp} H.~Wang \emph{et al.}, ``Unsupervised Representation Learning by
  Predicting Random Distances,'' in \emph{Proc. IJCAI}, 2020.
\bibitem{adbench} S.~Han \emph{et al.}, ``ADBench: Anomaly Detection Benchmark,'' in
  \emph{Proc. NeurIPS}, 2022.
\bibitem{odds} S.~Rayana, ``ODDS Library,'' Stony Brook University, 2016. [Online].
  Available: \url{https://odds.cs.stonybrook.edu}
\end{thebibliography}

\end{document}
